\section{Geographic Voronoi Districts}
\label{sec:phase2}

\subsection{Motivation: When Graph Distance Diverges from Geography}

Phase~1 \cvd{} measures proximity via BFS hop count on the census-tract
adjacency graph.
BFS hop count is a faithful proxy for geographic distance in most of the
contiguous United States: inland states with roughly regular census-tract
shapes produce adjacency graphs where hop count is proportional to road
distance, and Voronoi regions based on hop count are reasonable geographic
approximations.

However, BFS hop count diverges from geographic Euclidean distance in two
important cases:

\paragraph{Case 1: Coastal and waterfront states.}
States with substantial water boundaries (Florida, Louisiana, Michigan,
Maryland) have census tracts that are geographically far apart but connected
via many intermediate tracts along land corridors.
Two tracts on opposite shores of a bay may be 2 miles apart in Euclidean
distance but 50 hops apart in the adjacency graph (via the land route).
BFS Voronoi in such states assigns coastal tracts to ``wrong'' seeds based
on graph topology rather than actual geographic proximity, producing
elongated peninsular districts rather than compact coastal ones.

\paragraph{Case 2: Elongated rural states.}
States with elongated geographic shapes and irregular census-tract sizes
(Nevada, Montana, Wyoming) have rural tracts that are large in area but
contribute only 1 hop to BFS distance.
A rural tract covering 1{,}000 square miles is treated the same as an urban
tract covering 1 square mile in hop-count space.
BFS Voronoi places district boundaries in locations that are geometrically
natural in hop-count space but geographically arbitrary.

Geographic \cvd{} addresses both cases by replacing BFS hop count with
Euclidean distance on projected tract centroids.

\subsection{Geographic CVD: Algorithm Modification}

The geographic modification to Algorithm~\ref{alg:cvd} is minimal and is now
represented in the implementation by \texttt{VoronoiMetric::Geographic}:

\begin{enumerate}
\item \textbf{Data requirement}: Load tract centroid data into
  \texttt{LoadedGraph} as a vector of longitude/latitude pairs, one per tract.
  The loader now attempts to read centroid companion files and leaves the field
  empty when they are absent.  Geographic \cvd{} errors if the user requests
  it without centroid data.

\item \textbf{Metric change}: Replace $d_{\mathrm{BFS}}(t, \mathrm{seed}_j)$
  with $d_{\mathrm{Euc}}(t, \mathrm{seed}_j)$, the Euclidean distance between
  projected centroids.

\item \textbf{Centroid update}: Replace the graph-distance medoid with the
  population-weighted centroid:
  \[
    \mathrm{centre}_j = \frac{\sum_{t \in D_j} p_t \cdot \mathbf{x}_t}
                             {\sum_{t \in D_j} p_t},
  \]
  where $\mathbf{x}_t = (x_t, y_t)$ is the projected centroid of tract $t$.
  This is the true Euclidean centroid, not a medoid.
  The nearest tract to the centroid (by Euclidean distance) is used as the
  new seed for subsequent Voronoi assignments.

\item \textbf{Projection}: The current Rust implementation uses a simplified
  EPSG:5070-style Albers projection in \texttt{albers\_project}.  The
  specification identifies state-specific projections for Alaska, Hawaii, and
  Puerto Rico as future precision work; publication packages should record the
  projection used.
\end{enumerate}

\subsection{Expected Improvement over Phase 1}

Geographic \cvd{} is expected to yield a Polsby-Popper improvement over
graph-distance \cvd{} for states where BFS hop count diverges substantially
from geographic distance.  The earlier 5--10\% figure is a benchmark target,
not a publication claim until centroid-backed packages are archived.

The intuition: Euclidean Voronoi regions on a geographic plane are exactly
the geometric cells that minimise mean squared distance from seed to tract
centroid. These cells are circular in uniform-density regions, and elliptical
in anisotropic regions, both of which are more compact (higher PP) than the
irregular BFS Voronoi regions.
For states where hop count faithfully proxies geography (Iowa, Kansas),
the geographic improvement is expected to be small ($<3\%$).
For coastal states (Florida, Louisiana), the improvement is expected to be
larger in benchmark runs because geographic \cvd{} avoids the along-shore bias
of hop-count Voronoi.

\subsection{Implementation Status and Evidence Plan}

The implementation surface now exists:

\begin{itemize}
\item \texttt{LoadedGraph} has a \texttt{tract\_centroids} field.
\item \texttt{--cvd-metric geographic} dispatches to
  \texttt{split\_subgraph\_cvd\_geographic}.
\item \texttt{derive\_cvd\_geo\_seed} uses the distinct
  \texttt{"CVD\_GEO\_INIT\_"} prefix.
\item The CLI rejects geographic \cvd{} when centroid data are absent.
\end{itemize}

The remaining evidence plan is:

\begin{enumerate}
\item \textbf{Archive centroid-backed fixtures}:
  Commit at least one small synthetic centroid fixture and one real-data package
  where the centroid file, projection, metric, and seed derivation can be
  replayed.

\item \textbf{Validate on FL and LA}:
  Florida and Louisiana are the strongest tests because of their complex
  coastal geometries.
  Compare graph-distance and geographic \cvd{} PP on these states to confirm
  whether the expected improvement appears in real packages.

\item \textbf{Record projection metadata}:
  The run manifest should state the projection formula and whether the centroid
  data came from a companion file or from a geometry pipeline.
\end{enumerate}

\subsection{Connection to CVT Theory}

Geographic \cvd{} is directly the continuous CVT construction of
Du et al.~\citeyear{du1999} applied to the discrete redistricting domain.
With Euclidean distance, the geographic algorithm satisfies the two CVT conditions:
(1) each tract is assigned to its nearest seed (Voronoi condition);
(2) each seed is at the population-weighted centroid of its district
(centroidal condition, up to the discretisation of approximating a continuous
centroid by the nearest actual tract).

The continuous Lloyd guarantee applies before snapping the centroid back to a
real tract.  The implemented snap-to-nearest-tract step introduces a discrete
rounding error, so publication claims should state the weaker guarantee: the
algorithm is deterministic and replayable, and it follows the Lloyd update
heuristic, but monotone energy descent is not guaranteed at every snapped
iteration without an additional proof.

The convergence rate for geographic \cvd{} is expected to be faster than graph-distance \cvd{} in
Euclidean-like geographies: the Lloyd energy surface for Euclidean distance
is smoother and has fewer local minima than the graph-distance surface,
because Euclidean distance is a continuous metric while BFS hop count is
discrete and highly non-convex.
